% Propietario conservado: /Users/ruben/Documents/New project/output/EXTREMA_MEDIA_RAZON_RECIPROCIDAD_ES_EN_20260918/ARTICULO/ES/source/sections/electron.tex
\section{Structure algébrique de l'électron et échelle d'action}
\label{x_reciprocidad:sec:electron}

La construction électronique reçoit l'état arithmétique avec mémoire de trajectoire produit par APP--TRIT--TPK et ses coordonnées numériques. À ce stade du développement, \(\pi,\varphi,e\) et \(\alpha\) désignent les coordonnées déjà obtenues de la même généalogie, selon l'ordre constructif exposé précédemment. Ce ne sont pas des paramètres tirés d'une table physique. La section électronique est une réalisation ultérieure de cette structure discrète conjointe du continu : elle conserve la cellule d'origine, les deux feuillets arithmétiques, l'orientation, la retenue et la mémoire de retour avant de produire un scalaire.

Il convient de distinguer trois objets dès le départ. Le premier est la section
centrale et son histoire de transport ; le deuxième, la coordonnée
adimensionnelle qui en est obtenue ; le troisième, sa réalisation sur une
droite d'énergie ou de masse. L'égalité entre deux coordonnées numériques
n'identifie pas ces trois niveaux. En particulier, le nombre d'Euler \(e\),
la coordonnée électronique \(\varepsilon_e\) et une énergie de repos
\(E_e\) porteront des notations différentes.

L'objectif de cette section est d'établir la composition
\begin{equation}
 \varepsilon_e^{\beta}
 =\frac{\sqrt3}{4}
  \left[\exp\!\left(\frac{\varphi}{\pi^2}\right)-22\alpha^3\right]
  (1+15\Delta_4)
  R_{\mathrm{act}}^{\,80-54\alpha+6\Delta_4},
 \label{x_reciprocidad:eq:el-formula-completa}
\end{equation}
en expliquant séparément l'origine du préfacteur, les deux entiers de
l'étape basale, la normalisation de \(\Delta_4\), le rapport d'action et le
triplet de retour. La carte dimensionnelle ne s'ouvre qu'après
cette composition. Le résultat ne consiste pas à faire acquérir des unités à un nombre
sans dimension par sa seule ressemblance avec une grandeur
tabulée.

La composition réunit des facteurs aux fonctions différentes. Le tableau
suivant identifie leur signification et renvoie à la construction
correspondante ; les démonstrations se développent dans l'ordre de leurs
dépendances.

\begin{table}[htbp]
\centering
\small
\setlength{\tabcolsep}{4pt}
\renewcommand{\arraystretch}{1.18}
\begin{tabularx}{\linewidth}{@{}p{0.20\linewidth}Xp{0.24\linewidth}@{}}
\toprule
Facteur & Signification dans la composition & Construction et démonstration\\
\midrule
\(\sqrt3/4\)
 & Norme du commutateur des projecteurs dans le plan principal.
 & Équation~\eqref{x_reciprocidad:eq:el-prefactor}.\\
\(\exp(\varphi/\pi^2)\)
 & Caractère exponentiel de l'orientation électronique du quotient régional.
 & Équation~\eqref{x_reciprocidad:eq:el-bisagra} ; sections~\ref{x_reciprocidad:sec:medida-retorno-hojas}
   et~\ref{x_reciprocidad:sec:operador-retorno-marcado}.\\
\(-22\alpha^3\)
 & Contribution du complément régional au retour central.
 & Équations~\eqref{x_reciprocidad:eq:el-selector-enteros} et~\eqref{x_reciprocidad:eq:el-censo}.\\
\(1+15\Delta_4\)
 & Normalisation torsionnelle composée avec la multiplicité régionale
   et tritique.
 & Équations~\eqref{x_reciprocidad:eq:el-selector-enteros}, \eqref{x_reciprocidad:eq:el-delta}
   et~\eqref{x_reciprocidad:eq:el-basal}.\\
\(R_{\mathrm{act}}^{\kappa_e}\)
 & Quotient des sections d'action, élevé au caractère déterminé par les deux lecteurs du retour électronique.
 & Équations~\eqref{x_reciprocidad:eq:el-ratio}, \eqref{x_reciprocidad:eq:el-triple},
   \eqref{x_reciprocidad:eq:el-kappa} et~\eqref{x_reciprocidad:eq:el-beta}.\\
\bottomrule
\end{tabularx}
\caption{Facteurs de la composition électronique et provenance de leurs règles.}
\label{x_reciprocidad:tab:el-factores}
\end{table}


% BEGIN OWNER /Users/ruben/Documents/New project/output/EXTREMA_MEDIA_RAZON_RECIPROCIDAD_ES_EN_20260918/ARTICULO/ES/source/sections/electron_articulacion.tex
\paragraph{Identité de la section et composition de ses observables.}
\label{x_reciprocidad:el-ampl-articulacion}
L'identité structurelle de l'électron est introduite au moyen de la section
centrale et de ses opérations de transport. L'inversion de la carte APP
sélectionne la cellule $(5,5)$ par sa propriété de point fixe. Les deux feuillets
arithmétiques permettent d'étudier les déformations conservant la somme et
le résidu du produit. Le résultat distingue les deux orientations
$(8,2)$ et $(2,8)$ et conserve la différence d'une unité de quotient.
Tel est le premier niveau d'information : une cellule, une involution, une
fibre conservée et un défaut entier minimal. Aucune évaluation de masse
n'est nécessaire pour l'énoncer ou le démontrer.

Le passage aux projections arithmétiques introduit un deuxième niveau. Les
fibres de somme et de produit déterminent des espérances conditionnelles sur le
même espace fini. Leur commutateur mesure la dépendance à l'égard de
l'ordre de ces deux projections. La norme $\sqrt3/4$ résulte du spectre
de leur appariement ; le bloc qui atteint cette norme conserve en outre
les relations quadratiques démontrées ci-dessous. Le
préfacteur scalaire et l'algèbre locale ont donc une provenance commune, mais
contiennent des informations différentes. La norme retient la grandeur du
commutateur, tandis que les matrices conservent son orientation et sa loi de
composition.

Le transfert surfacique--volumique incorpore les coordonnées régionales
à un troisième niveau. Le calendrier des modes détermine la matrice
d'incidence et son action quadratique. La marque de clôture produit le rapport
$\pi/\varphi^2$ ; l'échange des deux arguments fournit l'
orientation $\varphi/\pi^2$ utilisée dans le caractère exponentiel
électronique. Les deux expressions sont reliées par une involution
explicite. Cette relation explique pourquoi une coordonnée régionale peut
apparaître dans l'équation électronique avec une puissance ou dans une position
différente de celle qu'elle occupait dans la première lecture : l'opération appliquée
fait partie de la donnée mathématique conservée.

Le terme cubique d'alpha appartient au complémentaire régional d'un retour
marqué. Le dénombrement de ses positions produit $27-5=22$, tandis que les
trois feuillets de chacune des cinq régions produisent $3\cdot5=15$.
Les démonstrations conservent le domaine de ces dénombrements et l'
orientation qui détermine leurs signes. La formule basale compose
donc une norme d'incompatibilité, un caractère de
transfert, un déficit de positions et une mémoire torsionnelle.
L'écriture sur une seule ligne exprime la composition de ces
opérations ; son explication exige de les reconstruire dans cet ordre.

La transition de la coordonnée basale à la coordonnée complète conserve
la même section centrale. Le facteur de retour utilise les registres
de l'histoire, et les deux opérateurs de lecture déclarés restituent le triplet
$(80,54,6)$ sur cette trajectoire. La multiplication par
$R_{\mathrm{act}}^{80-54\alpha+6\Delta_4}$ modifie l'évaluation,
mais maintient identifiés son origine et ses antécédents. La valeur
basale constitue une étape intermédiaire de la composition ; la valeur
complète inclut également cette contribution du retour. La comparaison
physique doit utiliser l'étape spécifiée par la formule évaluée.

L'échelle d'action intervient ensuite avec une autre fonction. Le rapport de
deux sections d'action est sans dimension et annule leur échelle commune.
La réalisation absolue utilise, en revanche, une base d'action, un
opérateur de lecture temporel et l'application vers l'énergie ou la masse. Cette distinction permet
de conserver simultanément la construction de l'échelle et l'identité
de la section électronique. Le quantum chronologique fournit une
référence d'échelle ; la section centrale détermine le facteur structurel
exprimé par rapport à elle. Les équations dimensionnelles de ce
chapitre maintiennent cette composition explicite.

La pertinence de l'électron dans l'argument général provient de cette
convergence d'opérations. La même arithmétique qui a produit les
régions détermine maintenant une cellule centrale, ses déformations et ses
projections. La constante de structure fine, construite auparavant, agit
dans cette composition comme coordonnée de clôture. L'information de
retour détermine l'évaluation complète, et la réalisation dimensionnelle
permet la comparaison ultérieure. Expliquer l'électron à partir de la structure
consiste à conserver cette succession d'objets, d'applications et de valeurs à chaque
étape de la démonstration.

% END OWNER


\subsection{Centre, orientation et vacance sur une fibre conservée}

Le support cellulaire considéré est \(\{1,\ldots,9\}^2\), avec les deux
évaluations APP, \(r+c\) et \(rc\). L'inversion centrale est
\begin{equation}
 \rho(r,c)=(10-r,10-c).
 \label{x_reciprocidad:eq:el-inversion-central}
\end{equation}
Le centre n'est pas choisi en fonction d'une masse : il se caractérise par une propriété du
support antérieure à sa lecture physique.

\begin{proposition}[Unicité du centre]
L'inversion \(\rho\) a un unique point fixe, \((5,5)\).
\end{proposition}
\begin{proof}
Les équations \(r=10-r\) et \(c=10-c\) équivalent à \(2r=2c=10\). Leur unique solution dans la carte nonadique est \((5,5)\).
\end{proof}

La propriété précédente identifie une cellule, et non une histoire profonde
unique. Une route ajoute le choix initial de l'orientation et ses
prolongations compatibles. Dans la carte utilisée ici, la section
orientée centrale commence en \((5,5)\) avec les deux orientations
positives ; elle sera notée \(\Gamma_e\). Le retour de sa projection
cellulaire n'efface ni les orientations transportées ni le registre de
retenue de la section avec sa mémoire de trajectoire.

La vacance est maintenant considérée, après la clôture de \(\alpha\), comme
une opération locale de cette section. Son arithmétique peut être isolée dans un
lemme fini ; cet isolement n'en fait pas un substitut du
générateur antérieur de constantes.

\begin{definition}[Fibre centrale et défaut de quotient]
La fibre additive centrale est
\[
 \mathcal F_{10}
 =\{(5+t,5-t):t\in\mathbb Z,\ |t|\leq4\}.
\]
Dans la carte \(n=\rho_9(n)+9q_9(n)\), avec \(\rho_9(n)\in\{1,\ldots,9\}\), on appelle défaut de quotient par rapport au centre
\[
 v(t)=q_9(25)-q_9(25-t^2),
\]
lorsque les deux produits ont le même résidu nonadique.
\end{definition}

\begin{lemma}[Vacance centrale minimale]
Les seuls points de \(\mathcal F_{10}\) qui conservent le résidu
multiplicatif du centre sont
\[
 (5,5),\qquad (8,2),\qquad (2,8).
\]
Aux deux points non centraux, le défaut de quotient est exactement un.
\end{lemma}
\begin{proof}
La somme est constamment égale à dix, tandis que
\[
 (5+t)(5-t)=25-t^2.
\]
La conservation du résidu équivaut à \(t^2\equiv0\pmod9\), c'est-à-dire à \(3\mid t\). Dans l'intervalle entier \([-4,4]\), les solutions sont \(t=-3,0,3\). Pour les deux solutions non nulles,
\[
 25=7+9\cdot2,\qquad16=7+9\cdot1.
\]
Le résidu sept est conservé et une unité de quotient est perdue. Il n'y a
pas d'autre déplacement admissible non central dans ce domaine, de sorte que le
défaut positif est minimal. L'échange \(t\mapsto-t\) permute les
deux sorties et conserve le défaut.
\end{proof}

La vacance n'est pas ici un zéro numérique : c'est une différence entre deux
états qui ont la même lecture résiduelle et des quotients distincts. Si l'on
ne conservait que le résidu, cette différence disparaîtrait. L'
orientation distingue en outre les deux points de sortie. Ni le défaut de
quotient ni cette orientation ne s'identifient automatiquement à une charge
électrique ou à une masse ; leur réalisation physique requiert les lecteurs
correspondants.


% BEGIN OWNER /Users/ruben/Documents/New project/output/EXTREMA_MEDIA_RAZON_RECIPROCIDAD_ES_EN_20260918/ARTICULO/ES/source/sections/centro_electronico_registros.tex
\subsubsection{Coordonnées centrales, orientation et registres régionaux}
\label{x_reciprocidad:el-centro:registros}

La structure centrale relie position, évaluation et transport. Dans la
cellule \((5,5)\), les deux feuillets d'APP donnent
\begin{equation}
 \mathcal A(5,5)=\bigl((1,1),(7,2)\bigr),\qquad
 5+5=1+9\cdot1,\qquad5\cdot5=7+9\cdot2.
 \label{x_reciprocidad:el-centro:app}
\end{equation}
Les deux unités du feuillet additif sont, dans cet ordre, le résidu et le quotient.
Leur coïncidence numérique conserve deux fonctions arithmétiques distinctes.
Sur la fibre de somme \(10\), les deux lacunes orientées vérifient
\begin{equation}
 \mathcal A(8,2)=\mathcal A(2,8)=\bigl((1,1),(7,1)\bigr).
 \label{x_reciprocidad:el-centro:vacancias}
\end{equation}
Le feuillet additif complet et le résidu multiplicatif sont conservés ;
la coordonnée de quotient du produit diminue exactement de
\(1\). Le transport retient également laquelle des deux positions
conjuguées a été atteinte.


% BEGIN OWNER /Users/ruben/Documents/New project/output/EXTREMA_MEDIA_RAZON_RECIPROCIDAD_ES_EN_20260918/ARTICULO/ES/source/figures/centro_y_vacancias.tex
\begin{figure}[H]
\centering
\begin{tikzpicture}[x=.67cm,y=.67cm,font=\footnotesize,>=Latex]
 \draw[step=1,black!12,line width=.3pt] (1,1) grid (9,9);
 \draw[->,black!60] (.7,1)--(9.5,1) node[right] {$i$};
 \draw[->,black!60] (1,.7)--(1,9.5) node[above] {$j$};
 \foreach \k in {1,...,9}{
  \node[below=2pt] at (\k,1) {$\k$};
  \node[left=2pt] at (1,\k) {$\k$};
 }
 \draw[black!65] (1,9)--(9,1);
 \draw[<->,line width=.9pt] (2,8)--(8,2);
 \foreach \x/\y in {2/8,5/5,8/2}{
  \filldraw[fill=white,line width=.8pt] (\x,\y) circle (3pt);
 }
 \fill (5,5) circle (1.5pt);
 \node[above right=3pt,fill=white,inner sep=1pt] at (2,8) {$(2,8)$};
 \node[above right=3pt,fill=white,inner sep=1pt] at (5,5) {$(5,5)$};
 \node[above right=3pt,fill=white,inner sep=1pt] at (8,2) {$(8,2)$};
 \node[anchor=west,align=left] at (10.5,7.6)
   {$i+j=10$\\$R^+=1,\quad q^+=1$};
 \node[anchor=west,align=left] at (10.5,5.4)
   {centre : $ij=25$\\$R^\times=7,\quad q^\times=2$};
 \node[anchor=west,align=left] at (10.5,3.1)
   {lacunes : $ij=16$\\$R^\times=7,\quad q^\times=1$};
\end{tikzpicture}
\caption{Centre électronique et lacunes conjuguées dans la carte APP.
La diagonale est la fibre additive \(\mathcal F_{10}\) ; le centre
est le point fixe de \((i,j)\mapsto(10-i,10-j)\). Les deux
déplacements de module \(3\) conservent la somme et le résidu
produit et réduisent d'une unité le quotient multiplicatif. Les
flèches représentent les déformations sur la fibre.}
\label{x_reciprocidad:el-centro:figura}
\end{figure}

% END OWNER


La symétrie centrale caractérise également sa persistance sous les changements de
carte. Si une transformation admissible \(f\) de l'atlas entrelace
l'inversion, \(f\rho=\rho f\), alors
\[
 \rho f(5,5)=f\rho(5,5)=f(5,5).
\]
L’unicité du point fixe impose \(f(5,5)=(5,5)\). C’est la
propriété interne qui rend stable la localisation centrale : toute
transformation de ce type conserve la cellule avant d’évaluer son
caractère électronique (section~\ref{x_reciprocidad:el-centro:registros}).

Le bloc matriciel central appartient à l'évaluation sur les deux
positions adjacentes \(4,5\). Sa décomposition, établie dans le
développement local du noyau, est
\[
 B_c=\begin{pmatrix}7&2\\2&7\end{pmatrix}=9P_++5P_-.
\]
Ici, \(9\) et \(5\) sont les valeurs propres de deux modes, tandis que
\((5,5)\) est une position de l'atlas. Les concaténations de ses
entrées donnent \(72+27=77+22=99\) et
\(77-22=55=54+1\). La première égalité conserve la somme totale sous
deux dispositions ; la seconde exprime le défaut entre les lectures
diagonale et antidiagonale. Ces opérations précèdent la composition
électronique et conservent leur provenance matricielle.

L'orientation initiale est une autre donnée. Dans la représentation du
retour électronique, on fixe \((h_0,v_0)=(1,1)\), c'est-à-dire une avance
positive dans les deux coordonnées. L'inversion simultanée tous les \(27\)
pas, lue par fenêtres de \(9\), détermine
\[
 \tau_e=(+1,+1,+1,-1,-1,-1,+1,+1,+1,-1,-1,-1).
\]
Les segments \((1,1,1)\) appartiennent à cette suite d'orientations :
chaque segment comprend trois fenêtres de même sens. La récurrence
de la section consacrée aux opérateurs de lecture démontre ce registre dans
\eqref{x_reciprocidad:eq:el-tau}, et ses deux évaluations restituent le triplet
\((80,54,6)\) qui intervient dans
\eqref{x_reciprocidad:eq:el-formula-completa}. La coordonnée de position, les
quotients APP et la suite des signes sont ainsi reliés par
des opérations explicites.

La signature \(555555\) appartient à une lecture régionale d'un autre
domaine. Pour une émission décimale \(U=(U_1,\ldots,U_6)\), on écrit
\(U_j=100d_{1j}+10d_{2j}+d_{3j}\), avec trois chiffres, y compris les
zéros initiaux. L'opérateur de lecture de la signature multiplicative est
\[
 \mathcal E_\times(U)
   =\bigl([d_{2j}-d_{1j}]_{10}\bigr)_{j=1}^{6}.
\]
Dans la région transversale
\(U_\pi^\perp=(501,614,498,169,272,272)\), les différences sont
\((-5,-5,5,5,5,5)\). Par conséquent,
\begin{equation}
 \mathcal E_\times(U_\pi^\perp)=(5,5,5,5,5,5),\qquad
 \Psi(U_\pi^\perp)=010211.
 \label{x_reciprocidad:el-centro:firma-transversal}
\end{equation}
La première application conserve une signature multiplicative du catalogue ;
la seconde conserve le mot ternaire commun aux cinq régions.
La figure \ref{x_reciprocidad:pi-estrella:figura} montre la position transversale de
cette émission et sa multiplicité \(432\).

\paragraph{Mémoire nécessaire au transport de la signature constante.}
La signature \(555555\) fournit également un exemple vérifiable de
l'information de trajectoire qu'exige TPK. Soit
\(\mathcal X_\times=D_9^2\times\{N,E,S,O\}\), avec \(324\)
positions orientées du canal produit. L'avance \(P\) déplace
d'une case dans la direction mémorisée et le demi-tour \(H\)
inverse cette direction. Dans cette représentation, l'opérateur de lecture
\(E_{\rm sig}\) somme par fenêtres les exposants discrets
\[
 \ell_9(1,2,4,8,7,5)=(0,1,2,3,4,5),\qquad
 \ell_9(3)=\ell_9(6)=\ell_9(9)=0.
\]
Pendant chacune des six fenêtres de neuf pas, on évalue
\(\ell_9(\rho_9(ij))\) avant chaque avance produit, sélectionnée
par TRIT aux instants \(2,5,8\pmod9\). La somme de la fenêtre est
réduite modulo \(10\). Après les pas \(27\) et \(54\), on inverse
la direction. Cette règle détermine les six composantes de
\(E_{\rm sig}\) à partir de la position et de l'orientation initiales.

L'énumération de ce domaine produit \(26\) signatures. La fibre de
\(555555\) contient \(24\) positions orientées ; après une
avance, leurs images se répartissent en trois signatures, avec huit
représentants chacune :
\[
 555177,\qquad555717,\qquad555771.
\]
Trois témoins, dans les coordonnées de \(1\) à \(9\), permettent de vérifier
directement la distinction :
\[
\begin{array}{c|c|c}
x&E_{\rm sig}(x)&E_{\rm sig}(Px)\\\hline
(3,5,N)&555555&555177\\
(3,4,N)&555555&555717\\
(3,4,S)&555555&555771
\end{array}
\]
La valeur commune de la signature ne détermine donc pas sa valeur suivante :
la classe de transport à l'intérieur de cette fibre fait partie de la donnée
opératoire. La partition stable minimale se calcule à partir de
l'équivalence des signatures \(\sim_0\), en raffinant par
\[
 x\sim_{n+1}y\quad\Longleftrightarrow\quad
 x\sim_n y,\quad Px\sim_n Py,\quad Hx\sim_n Hy.
\]
Le procédé s’achève parce que le domaine est fini. Toute congruence
stable qui raffine \(\sim_0\) raffine chaque \(\sim_n\), par récurrence ;
son point fixe est donc la congruence stable minimale requise.
Le recensement donne \(26\to28\to28\), avec l’histogramme
\(27\cdot8+108=324\). La seule signature qui se décompose est
\(555555\), en trois classes de huit éléments. Le certificat joint
reproduit l’énumération, les trois classes et la stabilité par les deux
opérateurs (section~\ref{x_reciprocidad:el-centro:registros}).

La multiplicité \(24\) appartient ici à la fibre du canal produit ;
la multiplicité \(432\) appartient à l'émission régionale complète
\(U_\pi^\perp\). La première démontre pourquoi la signature exige une mémoire
de transport ; la seconde situe cette signature dans le catalogue conjoint
des deux feuillets. Les deux lectures participent de la même architecture,
avec des domaines et des opérations spécifiés.

Cette distinction des domaines permet de réunir les occurrences de l'unité
et du centre avec leurs significations respectives :
\begin{center}
\small
\begin{tabularx}{\textwidth}{@{}lX@{}}
\toprule
Registro & Objet et fonction \\
\midrule
\((5,5)\) & Cellule fixe de l'inversion centrale de l'atlas APP.\\
\((R^+,q^+)=(1,1)\) & Résidu et quotient de l'évaluation additive centrale.\\
\((h_0,v_0)=(1,1)\) & Orientation initiale des deux déplacements.\\
\((+1,+1,+1)\) & Trois fenêtres consécutives de \(\tau_e\) d'orientation positive.\\
\(c=111111\) & Donnée de la cinquième frontière nonadique, section \ref{x_reciprocidad:mono-recta:desarrollo}.\\
\(\mathcal E_\times=555555\) & Signature de la région transversale de \(\pi\), à six positions.\\
\bottomrule
\end{tabularx}
\end{center}

La coordination électronique s'obtient en composant ces niveaux dans leur
ordre de dépendance. APP détermine le centre et les évaluations ;
TRIT sélectionne régimes et orientations ; TPK prolonge la trajectoire
et conserve ses registres. Les régions préalablement engendrées apportent
les caractères \(\pi,\varphi,e\) ; le registre dodécaphasique apporte
\(\alpha\) ; les opérateurs de lecture électroniques déterminent les facteurs de
l'équation \eqref{x_reciprocidad:eq:el-formula-completa}. Cette composition conserve
l'identité de la section à travers ses différentes
représentations, jusqu'à l'évaluation dimensionnelle et à la comparaison
métrologique ultérieure.

% END OWNER


\subsection{Norme du commutateur des projections arithmétiques}

On obtient le préfacteur électronique sans utiliser \(\alpha\), la torsion ou une unité dimensionnelle. Le domaine est la troisième réalisation d'APP, \(\mathcal A_3=\{1,\ldots,9\}^3\), munie de la mesure uniforme. Nous définissons
\[
 \Sigma(x)=\rho_9(x_1+x_2+x_3),\qquad
 \Pi(x)=\rho_9(x_1x_2x_3).
\]
Les espérances conditionnelles \(P_\Sigma\) et \(P_\Pi\) sont les
projections orthogonales sur les fonctions constantes sur chaque fibre
additive et multiplicative, respectivement. Ainsi, pour une fonction \(f\),
\[
 (P_\Sigma f)(x)
 =\frac{1}{|\Sigma^{-1}(\Sigma(x))|}
   \sum_{y:\,\Sigma(y)=\Sigma(x)}f(y),
\]
et de même pour \(P_\Pi\). Les deux projections agissent sur le même espace ; elles ne représentent pas deux supports arithmétiques indépendants.

\begin{lemma}[Deux projections et angles principaux]
Si \(s\in[0,1]\) est une valeur singulière de l'appariement entre des bases
orthonormales des images de deux projections \(P,Q\), son bloc non
trivial contribue à \(\|[P,Q]\|\) par
\(s\sqrt{1-s^2}\).
\end{lemma}
\begin{proof}
Dans le plan principal correspondant, on peut écrire
\[
 P=\begin{pmatrix}1&0\\0&0\end{pmatrix},\qquad
 Q=\begin{pmatrix}
 s^2&s\sqrt{1-s^2}\\
 s\sqrt{1-s^2}&1-s^2
 \end{pmatrix}.
\]
Leur commutateur est
\[
 [P,Q]=s\sqrt{1-s^2}
 \begin{pmatrix}0&1\\-1&0\end{pmatrix}.
\]
La matrice finale est orthogonale, donc sa norme est un. Les blocs communs ou orthogonaux ont un commutateur nul. La décomposition orthogonale en plans principaux donne l'affirmation et réduit la norme totale au maximum de ces contributions.
\end{proof}

\begin{proposition}[Norme exacte d'incompatibilité]
Sur \(\ell^2(\mathcal A_3)\),
\begin{equation}
 \|[P_\Sigma,P_\Pi]\|=\frac{\sqrt3}{4}.
 \label{x_reciprocidad:eq:el-prefactor}
\end{equation}
\end{proposition}
\begin{proof}
Soit \(N_{ab}=|\Sigma^{-1}(a)\cap\Pi^{-1}(b)|\). L'énumération entière
des trois symboles fournit
\[
 N=9\begin{pmatrix}
 0&1&1&0&1&1&0&1&4\\
 1&0&1&1&0&1&1&0&4\\
 1&0&2&0&1&2&0&0&3\\
 0&1&1&0&1&1&0&1&4\\
 1&0&1&1&0&1&1&0&4\\
 0&0&2&1&0&2&0&1&3\\
 0&1&1&0&1&1&0&1&4\\
 1&0&1&1&0&1&1&0&4\\
 0&1&2&0&0&2&1&0&3
 \end{pmatrix}.
\]
Chaque ligne a pour somme \(81\) ; les sommes des colonnes sont
\[
 (m_1,\ldots,m_9)=(36,36,108,36,36,108,36,36,297).
\]
Les indicatrices normalisées des fibres ont pour matrice
d'appariement \(C_{ab}=N_{ab}/\sqrt{81m_b}\). Par conséquent,
\[
 M=CC^{\mathsf T},\qquad
 M_{ac}=\sum_{b=1}^9\frac{N_{ab}N_{cb}}{81m_b}
\]
est une matrice rationnelle dont le spectre se compose des carrés des
valeurs singulières requises.

La multiplication exacte montre que \((1,\ldots,1)\),
\((1,-1,0)^3\) et \((1,1,-2)^3\) sont des vecteurs propres de valeurs propres
\(1\), \(1/4\) et \(7/132\), respectivement. Le sous-espace des
vecteurs à support dans \(\{3,6,9\}\) dont la somme est nulle a pour dimension
deux et pour valeur propre \(1/18\). Les vecteurs de somme nulle à support dans
\(\{1,4,7\}\) ou dans \(\{2,5,8\}\) forment un noyau de dimension
quatre. Ces sous-espaces sont indépendants et la somme de leurs dimensions vaut
neuf. Le spectre complet est donc démontré
\[
 \operatorname{spec}(M)=
 \left\{1,\frac14,\frac7{132},\frac1{18},\frac1{18},0,0,0,0\right\}.
\]
Le maximum de \(\lambda(1-\lambda)\) sur cet ensemble est \(3/16\), atteint en \(\lambda=1/4\). Le lemme précédent donne \(\sqrt{3/16}=\sqrt3/4\).
\end{proof}

Le mode \((1,-1,0)^3\) est précisément la distribution des valeurs du
sélecteur tritique de phase dans la carte nonadique. La norme scalaire retient
une mesure d'incompatibilité, mais ne conserve pas à elle seule
l'orientation de ce mode ni ne reconstruit les deux projecteurs. Cette mémoire
demeure dans l'état de route qui accompagne la lecture scalaire.

\begin{proposition}[Algèbre du bloc électronique local]
Dans le plan principal qui atteint \eqref{x_reciprocidad:eq:el-prefactor}, soient
\[
 P=\begin{pmatrix}1&0\\0&0\end{pmatrix},\qquad
 Q=\begin{pmatrix}1/4&\sqrt3/4\\\sqrt3/4&3/4\end{pmatrix},\qquad
 S=2P-I,\qquad J=\frac4{\sqrt3}[P,Q].
\]
Alors \(S^2=I\), \(J^2=-I\), \(SJ=-JS\), et l'algèbre réelle engendrée par \(S,J\) est \(M_2(\mathbb R)\), réalisation de \(\mathrm{Cl}_{1,1}\). En outre,
\[
 n_\pm=\frac{S\pm J}{2},\qquad
 n_\pm^2=0,\qquad n_+n_-+n_-n_+=I.
\]
\end{proposition}
\begin{proof}
On a \(S=\operatorname{diag}(1,-1)\) et
\(J=\bigl(\begin{smallmatrix}0&1\\-1&0\end{smallmatrix}\bigr)\).
La multiplication vérifie les trois relations. Les matrices \(I,S,J,SJ\) sont linéairement indépendantes et forment une base de \(M_2(\mathbb R)\), ce qui prouve l'affirmation sur l'algèbre. Enfin, \((S\pm J)^2=S^2+J^2\pm(SJ+JS)=0\), et la somme des produits croisés est \((S^2-J^2)/2=I\).
\end{proof}

Dans ce bloc coexistent une direction elliptique, une direction hyperbolique et deux
directions nulles paraboliques. Le résultat décrit une représentation
matricielle locale de l'arithmétique des deux feuillets. Il n'identifie pas toute APP
à \(M_2(\mathbb R)\), ni à un carré latin quantique ; de telles
affirmations exigeraient d'autres domaines et d'autres relations. L'exponentielle
n'appartient pas non plus exclusivement au régime parabolique : elle peut
s'appliquer à n'importe lequel des générateurs de ce bloc.


% BEGIN OWNER /Users/ruben/Documents/New project/output/EXTREMA_MEDIA_RAZON_RECIPROCIDAD_ES_EN_20260918/ARTICULO/ES/source/sections/helicidad_electron.tex
% Adición focal propuesta; no sustituye el álgebra electrónica existente.
% Inserción: después de la proposición «Álgebra del bloque electrónico local»
% y su párrafo explicativo, antes de «Transferencia entre hojas...».
% Procedencia y alcance: LOCALIZADORES_Y_ALCANCE.md, en esta carpeta.
\subsection{Représentation spinorielle du bloc électronique et hélicité}
\label{x_reciprocidad:sec:el-representacion-espinorial}

La structure électronique possède une représentation qui conserve davantage
d'information que son évaluation scalaire. Le plan principal des deux
projecteurs d'APP tridimensionnelle, sélectionné par le mode tritique
\((1,-1,0)^3\), a déjà fourni les opérateurs \(S\) et \(J\).
Le premier est une involution symétrique ; le second est le commutateur
normalisé et conserve l'orientation des feuillets. Leur composition permet
de construire explicitement une représentation des rotations et de séparer ses
deux composantes d'hélicité. Cette construction agit sur le bloc
électronique ; le registre TPK de la route et sa mémoire demeurent les données
de provenance du bloc.

Soit \(E_e\) le plan réel principal, muni du produit scalaire hérité
des espérances conditionnelles, et soit
\(\mathscr S_e=E_e\otimes_{\mathbb R}\mathbb C\) sa complexification.
Les opérateurs déjà obtenus satisfont
\[
 S^*=S,\qquad J^*=-J,\qquad
 S^2=I,\qquad J^2=-I,\qquad SJ=-JS.
\]
La complexification étend le corps de représentation. Elle conserve les
opérateurs réels et permet de distinguer l'unité scalaire \(i\) de
l'endomorphisme \(J\), bien que tous deux aient pour carré \(-1\) dans leurs
domaines respectifs.

\begin{proposition}[Triplet hermitien induit par les deux feuillets]
\label{x_reciprocidad:prop:el-terna-hermitica}
Les opérateurs
\begin{equation}
 \sigma_1=SJ,\qquad \sigma_2=-iJ,\qquad \sigma_3=S
 \label{x_reciprocidad:eq:el-terna-hermitica}
\end{equation}
sont hermitiens et de trace nulle. Ils satisfont
\begin{equation}
 \sigma_a\sigma_b
 =\delta_{ab}I+i\sum_{c=1}^3\epsilon_{abc}\sigma_c,
 \qquad
 \frac12\operatorname{tr}(\sigma_a\sigma_b)=\delta_{ab},
 \label{x_reciprocidad:eq:el-algebra-terna}
\end{equation}
où \(\epsilon_{123}=1\). Ils constituent donc une base orthonormée
de l'espace réel \(V_e\) des endomorphismes hermitiens de trace nulle de
\(\mathscr S_e\), pour le produit
\(\langle X,Y\rangle=\frac12\operatorname{tr}(XY)\).
\end{proposition}
\begin{proof}
Les relations précédentes donnent \((SJ)^*=(-J)S=SJ\) et
\((-iJ)^*=-iJ\). Les carrés des trois matrices valent \(I\).
De plus,
\[
 (SJ)(-iJ)=iS,\qquad (-iJ)S=iSJ,\qquad S(SJ)=J=i(-iJ).
\]
Inverser l'ordre change le signe. Cela établit la première identité
de \eqref{x_reciprocidad:eq:el-algebra-terna}. Dans la base orthonormée du plan principal,
les matrices sont
\[
 \sigma_1=\begin{pmatrix}0&1\\1&0\end{pmatrix},\quad
 \sigma_2=\begin{pmatrix}0&-i\\i&0\end{pmatrix},\quad
 \sigma_3=\begin{pmatrix}1&0\\0&-1\end{pmatrix}.
\]
Leur trace est nulle, et prendre les traces dans l'identité de produit établit
l'orthonormalité. Tout endomorphisme hermitien de trace nulle a la forme
\(\bigl(\begin{smallmatrix}z&x-iy\\x+iy&-z\end{smallmatrix}\bigr)\),
avec \(x,y,z\) réels ; il appartient donc à leur espace vectoriel réel engendré.
\end{proof}

Dans la représentation matricielle usuelle, ce triplet porte le nom de
matrices de Pauli. Son utilisation ici est postérieure à la construction de
\(P,Q,S,J\) : il s'obtient par les produits
\eqref{x_reciprocidad:eq:el-terna-hermitica}. L'espace des directions de cette
représentation est la sphère unité de \(V_e\), dont la forme positive
vient d'être construite au moyen de la trace.

\begin{proposition}[Générateurs de rotation et décomposition directionnelle]
\label{x_reciprocidad:prop:el-helicidad}
Définissons \(L_a=\sigma_a/2\). Alors
\begin{equation}
 [L_a,L_b]=i\sum_c\epsilon_{abc}L_c,
 \qquad \sum_{a=1}^3L_a^2=\frac34 I.
 \label{x_reciprocidad:eq:el-generadores-spin}
\end{equation}
Pour \(n=(n_1,n_2,n_3)\) avec \(\sum_an_a^2=1\), posons
\[
 H_e(n)=\sum_an_a\sigma_a,\qquad
 h_e(n)=\frac12H_e(n),\qquad
 \Pi_\pm(n)=\frac12\bigl(I\pm H_e(n)\bigr).
\]
La formule linéaire \(H_e(v)=\sum_av_a\sigma_a\) est également utilisée
pour \(v\) de norme arbitraire.
On a
\begin{equation}
 \mathscr S_e=\operatorname{im}\Pi_+(n)
             \mathbin{\oplus}\operatorname{im}\Pi_-(n),
 \qquad
 h_e(n)\Pi_\pm(n)=\pm\frac12\Pi_\pm(n).
 \label{x_reciprocidad:eq:el-descomposicion-helicidad}
\end{equation}
Les deux composantes sont des droites complexes orthogonales. Le transport de
rotation autour de \(n\) est
\begin{equation}
 U_n(\theta)=\exp\!\left(-\frac{i\theta}{2}H_e(n)\right)
 =\cos\frac\theta2\,I-i\sin\frac\theta2\,H_e(n),
 \qquad
 U_n(2\pi)=-I,\quad U_n(4\pi)=I.
 \label{x_reciprocidad:eq:el-retorno-spin}
\end{equation}
\end{proposition}
\begin{proof}
L'identité de produit du triplet donne les commutateurs de
\eqref{x_reciprocidad:eq:el-generadores-spin} ; la somme de leurs carrés vaut \(3I/4\).
Les termes antisymétriques s'annulent lorsqu'on multiplie
\(H_e(n)\) par lui-même, d'où \(H_e(n)^2=I\).
On en déduit
\[
 \Pi_\pm(n)^2=\Pi_\pm(n),\qquad
 \Pi_+(n)\Pi_-(n)=0,\qquad \Pi_+(n)+\Pi_-(n)=I.
\]
Ils sont hermitiens et de trace \(1\) ; leur rang vaut donc \(1\).
Multiplier par \(H_e(n)/2\) établit
\eqref{x_reciprocidad:eq:el-descomposicion-helicidad}. La série exponentielle se
sépare en puissances paires et impaires grâce à \(H_e(n)^2=I\), ce qui
produit \eqref{x_reciprocidad:eq:el-retorno-spin}. Elle établit également directement
l'unitarité et la loi \(U_n(\theta+\psi)=U_n(\theta)U_n(\psi)\).

La normalisation angulaire se vérifie sur \(V_e\). Pour tout
\(v\in\mathbb R^3\), la même identité de produit donne
\[
 U_n(\theta)H_e(v)U_n(\theta)^*
 =H_e\!\left(v\cos\theta+(n\mathbin{\times}v)\sin\theta
       +n(n\mathbin{\cdot}v)(1-\cos\theta)\right).
\]
L'action induite est une rotation d'angle \(\theta\) et d'axe \(n\).
Le facteur \(1/2\) dans le générateur spinoriel est donc celui qui
correspond à cette paramétrisation angulaire. Une rotation complète rétablit
la direction dans \(V_e\) et change le signe du vecteur de
\(\mathscr S_e\) ; deux rotations complètes rétablissent les deux.
\end{proof}

La représentation est irréductible : un sous-espace complexe invariant sous
les trois matrices serait invariant sous toute l'algèbre \(M_2(\mathbb C)\),
qu'elles engendrent. Ses seuls sous-espaces invariants sont \(0\) et
\(\mathscr S_e\). Les relations
\eqref{x_reciprocidad:eq:el-generadores-spin} identifient ainsi la représentation de
spin \(1/2\). Si \(n\) est la direction de propagation dans une
réalisation cinématique, \(h_e(n)\) est son hélicité sans dimension, de
valeurs \(\pm1/2\). La direction cinématique est déclarée lors de cette
réalisation ; les deux projecteurs sont définis pour chaque direction
avant la sélection d'un état de propagation particulier.

\paragraph{Conservation de la structure électronique et types d'inversion.}
L'inversion de direction possède une loi exacte :
\[
 H_e(-n)=-H_e(n),\qquad \Pi_\pm(-n)=\Pi_\mp(n).
\]
Cette opération échange les étiquettes d'hélicité relatives à une
direction. La conjugaison de charge agit, en revanche, sur la signature
orientée de la route matérielle ; l'inversion cellulaire agit sur les
coordonnées de l'atlas. Leurs domaines et leurs actions demeurent distincts.
L'unicité de la cellule \((5,5)\) est compatible avec les deux droites de
\eqref{x_reciprocidad:eq:el-descomposicion-helicidad} : une cellule de l'atlas n'est pas un
vecteur de la représentation spinorielle qu'elle porte.

L'évaluation électronique \(\varepsilon_e^\beta\) conserve sa
formule et ses opérateurs de lecture. Sur \(\mathscr S_e\), sa réalisation scalaire
est \(\varepsilon_e^\beta I\), qui commute avec les générateurs et avec
\(\Pi_\pm(n)\). La décomposition d'hélicité ajoute l'information
rotationnelle sans modifier l'évaluation ni convertir le spin en un
coefficient multiplicatif de la masse.

La chiralité correspond à la graduation de la représentation relativiste
utilisée : une involution \(\Gamma_{\rm ch}\) détermine ses
projecteurs \((I\mp\Gamma_{\rm ch})/2\). L'hélicité de
\eqref{x_reciprocidad:eq:el-descomposicion-helicidad} dépend en outre de \(n\).
L'holonomie d'un fibré réel de Möbius décrit, quant à elle, le changement
de signe d'une droite lorsqu'on parcourt un lacet de sa base. La loi de cette droite,
le retour spinoriel \eqref{x_reciprocidad:eq:el-retorno-spin} et l'holonomie nonadique
de phase avec avance de mémoire sont des lois de transport sur des objets
spécifiés. Leur comparaison conserve ces objets : la forme hélicoïdale
d'une trajectoire ne remplace pas la représentation qui vient d'être
construite, ni ne fixe à elle seule la chiralité ou la charge.

% END OWNER


\subsection{Transfert entre feuillets et coefficients de l'étape basale}

Le calendrier tritique primitif visite les modes
\((\Sigma,\Pi,\Pi)\). En comptant les couples consécutifs, y compris la
fermeture du bloc, on obtient une transition \(\Sigma\to\Pi\), une
\(\Pi\to\Pi\) et une \(\Pi\to\Sigma\). Leur incidence primitive est
\[
 F_{\mathrm{av}}=\begin{pmatrix}0&1\\1&1\end{pmatrix}.
\]
Cette réalisation matricielle s'obtient après le calendrier. Elle n'est pas employée pour remplacer le lecteur coinductif qui a déjà généré \(\varphi\). Son polynôme caractéristique est \(X^2-X-1\), et la reconnaissance des deux racines comme \(\varphi\) et \(-\varphi^{-1}\) décrit l'action de cette sortie dans le transfert modal. L'action quadratique a pour valeurs propres \(\varphi^2,-1,\varphi^{-2}\) ; la dernière est le mode stable positif. La marque régionale \(\pi\), appliquée une seule fois, donne \(\pi/\varphi^2\).

L'espace des trajectoires, sa mesure et l'opérateur de clôture sont
construits dans les paragraphes~\ref{x_reciprocidad:sec:medida-retorno-hojas}
et~\ref{x_reciprocidad:sec:operador-retorno-marcado}. La formule
\eqref{x_reciprocidad:eq:retorno-marcado-finito} conserve sa provenance comme quotient
de retours ; l'application électronique compose ensuite cette lecture
avec l'involution suivante. La signature régionale \(555555\) accompagne
la région de clôture, dont le transport conserve l'état à mémoire ;
le centre \((5,5)\) fixe la cellule d'origine de la section électronique
orientée. Les deux informations convergent dans la composition et conservent
leurs domaines respectifs.

\begin{lemma}[Échange des deux orientations]
Soient \(\mathcal R(x,y)=x/y^2\) et
\(\mathscr S_{\pi\varphi}(x,y)=(y,x)\), pour \(x,y>0\).
Alors
\begin{equation}
 \begin{aligned}
 \mathscr S_{\pi\varphi}^{\,2}&=I,\qquad
 \mathcal R(\pi,\varphi)=\frac{\pi}{\varphi^2},\\
 (\mathcal R\circ\mathscr S_{\pi\varphi})(\pi,\varphi)
 &=\frac{\varphi}{\pi^2}
 =\frac1{\varphi^3(\pi/\varphi^2)^2}.
 \end{aligned}
 \label{x_reciprocidad:eq:el-bisagra}
\end{equation}
\end{lemma}
\begin{proof}
L'échange appliqué deux fois est l'identité. De plus,
\(\varphi^3(\pi/\varphi^2)^2=\pi^2/\varphi\) ; prendre son inverse prouve
la dernière égalité.
\end{proof}

Le lecteur électronique utilise l'orientation \(\varphi/\pi^2\),
suivie du caractère exponentiel. L'identité précédente explique
l'échange, mais ne remplace pas la règle qui choisit ce caractère ni ne produit
à elle seule les entiers de la correction. Ceux-ci ont une autre provenance.

Le retour central compatible avec la phase a pour longueur minimale vingt-sept. Avec le registre de phase et de feuillet, son espace fini est
\[
 \mathcal K_e=\mathbb Z/9\mathbb Z\times\mathbb F_3.
\]
Les cinq régions de \(\pi\) sont des identités généalogiques distinctes,
et non cinq copies interchangeables d'un développement visible. Dans la jauge
régionale fixée par la construction, leur injection dans le retour possède
les emplacements
\begin{equation}
 (4,0),\quad(5,0),\quad(6,0),\quad(6,1),\quad(6,2).
 \label{x_reciprocidad:eq:el-cinco-lugares}
\end{equation}
L'orientation transversale, le retour muté et les trois régions positives restent séparés par ce registre. Les changements admissibles de jauge transportent les étiquettes conjointement ; ils ne sélectionnent pas à nouveau les positions à partir d'une valeur électronique.

\begin{proposition}[Déficit et mémoire retenue]
Une fois fixés le retour central et l'injection régionale précédente, le déficit
orienté de frontière et la mémoire torsionnelle sont
\begin{equation}
 n_e=-|\mathcal K_e\setminus\iota(\mathscr R_\pi)|=-22,
 \qquad
 m_e^{\mathrm{tor}}=|\mathscr R_\pi\times\mathbb F_3|=15.
 \label{x_reciprocidad:eq:el-selector-enteros}
\end{equation}
\end{proposition}
\begin{proof}
Les cinq emplacements de \eqref{x_reciprocidad:eq:el-cinco-lugares} sont distincts. Par
conséquent, \(|\mathcal K_e|=9\cdot3=27\) et le complémentaire de l'image
a pour cardinal \(27-5=22\). La convention d'orientation enregistre comme
négatif le déficit retiré de la frontière. La mémoire retenue possède
trois feuillets par région, de sorte que son cardinal positif est
\(5\cdot3=15\).
\end{proof}

Le même couple possède le contrôle entier de la carte APP de la famille :
\begin{equation}
 396=18+72+108+198,
 \qquad \frac{396}{18}=22,
 \qquad \frac{270}{18}=15=3\cdot5=\binom62.
 \label{x_reciprocidad:eq:el-censo}
\end{equation}
Le poids \(18\) appartient au noyau central \(2\times2\), et \(270\) au canal global déclaré. Ces égalités contrôlent la normalisation du décompte ; l'argument de retour fixe en outre les signes. Le quotient des anneaux alternés, \(396/(18+108)=22/7\), est rationnel et ne s'identifie pas à \(\pi\). De même, \(\binom62\) ne signifie pas \(6/2\) : il compte les paires non ordonnées d'un ensemble de six éléments.

L'évaluation du déficit dans la troisième réalisation symétrique utilise
\(\alpha^3\). Cette puissance est le produit de trois copies du scalaire
de clôture déjà obtenu ; elle n'affirme pas de règle universelle faisant
correspondre toute dimension à la même puissance de \(\alpha\).
Le domaine et le lecteur fixés sont essentiels à la composition.

\subsection{La torsion globale et la coordonnée basale}

\begin{definition}[Normalisation torsionnelle électronique]
On distingue le défaut réduit \(d_4\) et le lecteur complet
\(\Delta_4\) :
\begin{equation}
 d_4=\frac{\pi-e\log\pi}{270},\qquad
 \Delta_4=d_4\left(1+\frac{\pi}{729}\right).
 \label{x_reciprocidad:eq:el-delta}
\end{equation}
Le dénominateur \(270\) correspond au canal du recensement, tandis que le facteur \(1+\pi/729\) conserve le couplage avec la cellule cubique de \(729\) positions. La lecture électronique utilise \(\Delta_4\) aussi bien dans son étape basale que dans l'exposant de retour.
\end{definition}

Cette définition agit sur \(\pi\) et \(e\) préalablement générés. Le
logarithme est une opération ultérieure dans leur réalisation réelle, et non une
entrée rétrospective de valeurs conventionnelles. La torsion est globale :
elle ne se détermine pas indépendamment pour chaque section au moyen d'une masse
observée.

\begin{lemma}[Les deux normalisations ne sont pas interchangeables]
On a
\begin{equation}
 \Delta_4-d_4=\frac{\pi}{729}d_4.
 \label{x_reciprocidad:eq:el-delta-diferencia}
\end{equation}
En particulier, \(d_4>0\) et \(\Delta_4>d_4\) pour les coordonnées
considérées.
\end{lemma}
\begin{proof}
L'identité est distributive. Pour la positivité, la fonction
\(f(x)=x-e\log x\), sur \(x>0\), a pour dérivée \(1-e/x\) et
pour dérivée seconde \(e/x^2>0\). Son unique minimum est en \(x=e\), où
elle vaut zéro. Comme \(\pi\ne e\), on a \(f(\pi)>0\). Les deux facteurs
de \eqref{x_reciprocidad:eq:el-delta} sont positifs et \(1+\pi/729>1\).
\end{proof}

\begin{remark}[Conservation de la normalisation]
Supprimer le second facteur de \eqref{x_reciprocidad:eq:el-delta} change le lecteur ; ce n'est pas une simplification algébrique. Si l'on introduit \(d_4\) comme coordonnée auxiliaire, le même objet doit s'écrire \(\Delta_4=(1+\pi/729)d_4\) à chacune de ses occurrences. Maintenir les entiers \(15\) et \(6\) et remplacer seulement \(\Delta_4\) par \(d_4\) produit une autre formule.
\end{remark}

\begin{definition}[Coordonnée électronique basale]
La composition de l'incompatibilité, du transfert modal, du déficit et de la mémoire retenue
est
\begin{equation}
 \varepsilon_e^{(0)}
 =\frac{\sqrt3}{4}
  \left[\exp\!\left(\frac{\varphi}{\pi^2}\right)-22\alpha^3\right]
  (1+15\Delta_4).
 \label{x_reciprocidad:eq:el-basal}
\end{equation}
\end{definition}

L'équation réunit trois opérations de natures différentes. La norme
\(\sqrt3/4\) quantifie l'incompatibilité des projections des
feuillets ; son origine est un angle principal des fibres d'APP.
L'exponentielle agit sur l'orientation \(\varphi/\pi^2\) du
lecteur de transfert. La soustraction cubique enregistre les \(22\) positions
complémentaires du retour marqué, et le facteur torsionnel conserve les
\(15\) composantes correspondant à trois feuillets par région. La
composition conserve ces domaines même si son évaluation finale est scalaire.

La présence de \(\pi\), \(\varphi\) et \(e\) possède ainsi des significations
opératoires. Le nombre d'Euler intervient par l'exponentielle et par
la normalisation logarithmique de \(\Delta_4\) ; \(\varphi\) et \(\pi\)
interviennent dans le transfert orienté. La constante \(\alpha\)
agit après sa propre construction au moyen du terme cubique
de clôture. Le retour chronologique doit encore opérer sur cette coordonnée
basale. La séparation des étapes permet d'expliquer ce qui change lorsqu'on modifie
une orientation, une normalisation ou un registre, tout en conservant la structure
centrale qui identifie la section électronique.

Les parenthèses de cette équation font partie de son contenu. Le terme \(-22\alpha^3\) est soustrait \emph{en dehors} de l'exponentielle ; la mémoire \(1+15\Delta_4\) multiplie toute l'expression entre crochets. Introduire le déficit dans l'exposant ou lui appliquer une torsion différente altérerait la composition.

\begin{proposition}[Positivité basale]
Pour \(0<\alpha<10^{-2}\), \(\varphi>0\) et la normalisation \eqref{x_reciprocidad:eq:el-delta}, on a \(\varepsilon_e^{(0)}>0\).
\end{proposition}
\begin{proof}
L'exponentielle est supérieure à un, et
\(22\alpha^3<22\cdot10^{-6}<1\). Le crochet est donc positif.
On a également \(1+15\Delta_4>1\) et \(\sqrt3/4>0\).
\end{proof}

L'évaluation ultérieure de \eqref{x_reciprocidad:eq:el-basal} commence par
\[
 \varepsilon_e^{(0)}=0.5109989224880660725\ldots.
\]
Ce nombre est encore une coordonnée sans dimension et ne constitue pas la dernière étape : il reste à appliquer la mémoire de retour de la route elle-même.

\subsection{Échelle déterminantale d'action et valuation décimale}

L'échelle absolue et le correcteur relatif proviennent d'objets
distincts. La première utilise le bloc local du registre dodécaphasique
de pas trois,
\begin{equation}
 H_4=\begin{pmatrix}
 1&1&1&1\\1&1&-1&-1\\1&-1&1&-1\\1&-1&-1&1
 \end{pmatrix},\qquad
 G_H=\mathbb Z^4/H_4\mathbb Z^4.
 \label{x_reciprocidad:eq:el-h4}
\end{equation}
L'opération utilise un quotient local. Les trois blocs du registre
complet ne sont pas remplacés par une tensorisation qui multiplierait à
nouveau le nombre de feuillets de ce lecteur.

\begin{lemma}[Nombre de feuillets du quotient local]
La forme normale de Smith de \(H_4\) est
\(\operatorname{diag}(1,2,2,4)\). Par conséquent,
\[
 G_H\simeq\mathbb Z/2\mathbb Z\oplus\mathbb Z/2\mathbb Z
 \oplus\mathbb Z/4\mathbb Z,\qquad |G_H|=16.
\]
\end{lemma}
\begin{proof}
Le plus grand commun diviseur des entrées vaut un. Les mineurs d'ordre deux
sont pairs et l'un d'eux vaut \(-2\), de sorte que leur plus grand commun diviseur
vaut deux. Comme \(H_4H_4^{\mathsf T}=4I\) et \(\det H_4=-16\), les
cofacteurs valent \(\pm4\) ; le troisième diviseur déterminantal vaut quatre.
Le quatrième vaut \(16\). Les quotients successifs des diviseurs
\(1,2,4,16\) sont \(1,2,2,4\). Leur produit donne l'ordre du quotient.
\end{proof}

\begin{definition}[Lecteur déterminantal d'action]
La norme multiplicative de la section scalaire \(\alpha\), constante
sur les feuillets de \(G_H\), et une application de l'auto-échelle
\(\varphi\) produisent
\begin{equation}
 \Sigma_H=\varphi\,\mathrm N_{G_H}(\alpha),\qquad
 \mathrm N_{G_H}(\alpha)=\prod_{g\in G_H}\alpha=\alpha^{16}.
 \label{x_reciprocidad:eq:el-lector-determinantal}
\end{equation}
La carte décimale associée détermine
\(\operatorname{ord}_{10}(\Sigma_H)=\lfloor\log_{10}\Sigma_H\rfloor\).
\end{definition}

La puissance seize provient ainsi du nombre de feuillets du quotient, une
fois ce lecteur local fixé. Le facteur \(\varphi\) agit une fois sur
la base de la droite d'action. Le calcul de Smith, à lui seul, ne
sélectionnerait pas n'importe quelle fonctionnelle possible sur le quotient : la norme
multiplicative et l'auto-échelle font partie de l'opération déclarée.

\begin{theorem}[Décade de la section d'action]
Le lecteur \eqref{x_reciprocidad:eq:el-lector-determinantal}, appliqué aux coordonnées
HMT déjà engendrées, satisfait
\begin{equation}
 10^{-34}<\Sigma_H<10^{-33},\qquad
 \operatorname{ord}_{10}(\Sigma_H)=-34.
 \label{x_reciprocidad:eq:el-decada}
\end{equation}
\end{theorem}
\begin{proof}
Les intervalles rationnels suivants, beaucoup plus grossiers que les cylindres
disponibles, suffisent :
\[
 \frac{1618}{1000}<\varphi<\frac{1619}{1000},\qquad
 \frac{7297}{10^6}<\alpha<\frac{7298}{10^6}.
\]
La fonction \((x,a)\mapsto xa^{16}\) est croissante en ses deux arguments
positifs. Les deux inégalités entières
\[
 1618\cdot7297^{16}>10^{65},\qquad
 1619\cdot7298^{16}<10^{66}
\]
impliquent
\[
 10^{-34}
 <\frac{1618}{1000}\left(\frac{7297}{10^6}\right)^{16}
 <\Sigma_H
 <\frac{1619}{1000}\left(\frac{7298}{10^6}\right)^{16}
 <10^{-33}.
\]
La définition de l'ordre décimal conclut la preuve. Ni une
unité SI ni une masse de référence n'intervient dans ces inégalités.
\end{proof}

La section normalisée commence par
\(\Sigma_H=1.0462438381\ldots\times10^{-34}\). Sa mantisse ne
s'identifie pas à la coordonnée d'action avant ou après le retour.
Le lecteur déterminantal fixe l'ordre décimal ; les deux coordonnées de
retour s'obtiennent au moyen des lecteurs suivants.


\subsection{Sections d'action et quotient adimensionnel de retour}

Soit \(\mathcal L_S\) une droite orientée d'action et \(\mathcal U_S\)
sa base positive. La composition électronique utilise deux sections de
cette droite et leur quotient sans dimension. Sa construction requiert la
décade déterminantale déjà obtenue, le caractère d'action d'ordre cinq
et la contribution régionale du retour.

\subsubsection{Caractère d'action et provenance de ses coefficients}
\label{x_reciprocidad:sec:revision-planck}

La constante de Planck réduite, \(\hbar\), est l'objet physique qui
motive la réalisation de la section d'action angulaire. Sa position dans
la dépendance électronique doit être spécifiée à deux niveaux. La norme
déterminantale \(\Sigma_H\) détermine la décade de l'échelle. Le caractère
\(H_5\) et le retour régional déterminent ensuite les coordonnées des
sections d'action. Cette séparation permet d'exposer le résultat
nécessaire à l'électron sans anticiper les développements complets
d'autres secteurs de la Mécanique Dimensionnelle.

La provenance du caractère d'ordre cinq peut être explicitée au moyen
d'invariants déjà fixés dans le support. Dans la carte centrale, on considère
\[
 B_c=\begin{pmatrix}7&2\\2&7\end{pmatrix},
 \qquad \lambda_+=9,\quad\lambda_-=5.
\]
Le calendrier est de longueur \(12\), l'orbite du pas \(3\) est de
longueur \(4\), et le quotient de dénombrement pertinent est
\(104976/1944=54\). Les coefficients du caractère satisfont
\begin{equation}
 \frac9{16}=\frac{\lambda_+}{4^2},\quad
 \frac59=\frac{\lambda_-}{\lambda_+},\quad
 \frac7{48}=\frac{\operatorname{tr}(B_c)/2}{4\cdot12},\quad
 \frac1{54}=\frac{1944}{104976}.
 \label{x_reciprocidad:eq:revision-h5-coeficientes}
\end{equation}
L'attribution de ces invariants aux degrés \(2,3,4,5\), avec
l'alternance des signes de \eqref{x_reciprocidad:eq:el-h5}, appartient à la définition du
caractère analytique utilisé. L'expliciter est essentiel : l'ensemble des
cardinaux, considéré sans le caractère ni l'ordre de ses degrés,
admettrait d'autres expressions. Le contenu constructif réside dans la
composition du support, de l'orientation et de la règle analytique déclarée.

Sur cette base, le caractère de cinquième ordre et le terme régional de
retour ont pour expressions
\begin{align}
 H_5={}&\frac12\sqrt{\frac{1000\alpha}{\varphi}}-\alpha
   +\frac9{16}\alpha^2-\frac59\alpha^3
   +\frac7{48}\alpha^4-\frac1{54}\alpha^5,
   \label{x_reciprocidad:eq:el-h5}\\
 D_A={}&\exp\!\left(-\frac{100\pi\alpha}{9}\right),
   \label{x_reciprocidad:eq:el-da}\\
 \mathcal C_\pi(\alpha)={}&169\alpha^6+
       \frac{D_A\alpha^7}{1-D_A\alpha},
   \label{x_reciprocidad:eq:el-cpi}\\
 \eta_{\mathrm{ret}}={}&H_5-\frac{90}{\pi}\mathcal C_\pi(\alpha).
   \label{x_reciprocidad:eq:el-eta}
\end{align}
Ce sont des lectures ultérieures de la fermeture dodécaphasique et de son
incidence régionale. L'entier \(169\) est conservé comme le sélecteur
régional déjà construit, et les coefficients de \(H_5\) comme ceux de son
caractère local ; ils ne sont pas remplacés par une interpolation de la constante de
Planck. Les formules explicitent exactement les lectures nécessaires
à l'étape électronique. Il n'est pas nécessaire d'introduire ici d'autres
fonctionnelles de la théorie des masses.

Pour \(0<\alpha<1\), on a \(0<D_A<1\), de sorte que
\(D_A\alpha<1\). Le terme rationnel de \eqref{x_reciprocidad:eq:el-cpi} conserve la
somme complète de la queue géométrique
\begin{equation}
 \frac{D_A\alpha^7}{1-D_A\alpha}
 =\sum_{n=0}^{\infty}D_A^{n+1}\alpha^{n+7}.
 \label{x_reciprocidad:eq:el-cola-regional}
\end{equation}
Il ne s'agit pas d'une troncature que l'on pourrait modifier silencieusement en changeant la
précision. Après le terme d'indice \(N\), la queue restante est
\(D_A^{N+2}\alpha^{N+8}/(1-D_A\alpha)\), ce qui permet de contrôler
l'évaluation sans recourir à une grandeur cible.

\begin{definition}[Sections d'action et lecteur de retour]
Les sections avant et après le retour sont
\begin{equation}
 \hbar_{\mathrm{pre}}=H_5\,10^{-34}\mathcal U_S,
 \qquad
 \hbar_{\mathrm{ret}}=\eta_{\mathrm{ret}}\,10^{-34}\mathcal U_S.
 \label{x_reciprocidad:eq:el-secciones-accion}
\end{equation}
Lorsque \(\eta_{\mathrm{ret}}>0\), on définit
\begin{equation}
 \delta_{\mathrm{ret}}=H_5-\eta_{\mathrm{ret}},\qquad
 R_{\mathrm{act}}=\frac{\hbar_{\mathrm{pre}}}{\hbar_{\mathrm{ret}}}
 =\frac{H_5}{\eta_{\mathrm{ret}}}.
 \label{x_reciprocidad:eq:el-ratio}
\end{equation}
\end{definition}

\begin{proposition}[Positivité et invariance du rapport]
Dans les intervalles utilisés dans la preuve de \eqref{x_reciprocidad:eq:el-decada},
\(0<\eta_{\mathrm{ret}}<H_5\), et donc
\(R_{\mathrm{act}}>1\). Le rapport ne change pas lorsqu'on remplace la base
\(\mathcal U_S\) par toute autre base positive commune.
\end{proposition}
\begin{proof}
Les termes de \(\mathcal C_\pi\) sont positifs. Pour obtenir une borne
élémentaire, il suffit d'utiliser \(\alpha<0.008\), \(\varphi<1.619\),
\(\alpha>0.007\), \(D_A<1\) et \(\pi>3\). La racine dans
\eqref{x_reciprocidad:eq:el-h5} est supérieure à un, tandis que la somme des termes
négatifs est inférieure à \(0.008001\) ; en particulier, \(H_5>0.99\).
De même,
\[
 0<\frac{90}{\pi}\mathcal C_\pi
 <30\left(169(0.008)^6+
           \frac{(0.008)^7}{1-0.008}\right)<10^{-6}.
\]
Ainsi, \(\eta_{\mathrm{ret}}>0.989999>0\) et
\(\eta_{\mathrm{ret}}<H_5\). Enfin, le changement de base divise
les deux coordonnées d'action par le même facteur positif, qui se simplifie dans
leur quotient.
\end{proof}

L'évaluation ultérieure fournit
\begin{align*}
 H_5&=1.0545718176461544696\ldots,\\
 \eta_{\mathrm{ret}}&=1.0545718169150390611\ldots,\\
 R_{\mathrm{act}}&=1.0000000006932817630\ldots.
\end{align*}
La différence et le quotient sont deux lectures des mêmes sections,
l'une additive et l'autre multiplicative. Elles ne représentent pas deux constantes de
Planck indépendantes. Le passage de l'action angulaire à l'action par tour
complet s'exprime, dans chaque section, par
\(h=2\pi\hbar\). La carte ultérieure
\(\mathcal U_S\mapsto\mathrm{J\,s}\) attribue une unité, mais ne
rétroagit pas sur l'ordre décimal déjà démontré.

Le retour régional peut aussi s'écrire dans une carte angulaire.
Cette reparamétrisation n'est pas une dépendance supplémentaire indispensable :
\eqref{x_reciprocidad:eq:el-h5}--\eqref{x_reciprocidad:eq:el-eta} fournissent directement
\(H_5\), \(\eta_{\mathrm{ret}}\) et \(R_{\mathrm{act}}\). Si l'on
introduit l'angle et le contre-angle, tous deux doivent être transportés dans la même
carte ; degrés et radians ne sont pas mélangés dans une différence.


% BEGIN OWNER /Users/ruben/Documents/New project/output/EXTREMA_MEDIA_RAZON_RECIPROCIDAD_ES_EN_20260918/ARTICULO/ES/source/sections/revision_planck.tex
\subsubsection{Continuation régionale et équation de renouvellement}

Le sélecteur \(169\) détermine l'impulsion de degré six. Le prolongement
ultérieur s'organise selon le transfert \(D_A\) et une normalisation
unitaire de la droite déterminantale. Ces règles admettent une formulation
en séries formelles avant la substitution \(z=\alpha\), avec \(D_A\) déjà
construit. Écrivons
\[
 \mathscr C(z)=169z^6+R(z),\qquad R(z)\in z^7\mathbb R[[z]].
\]
La concaténation d'un prolongement augmente le degré d'une unité et
multiplie son poids par \(D_A\). Le premier prolongement strict a pour
poids \(D_A\). Son équation de renouvellement est donc
\begin{equation}
 R(z)=D_Az^7+D_AzR(z).
 \label{x_reciprocidad:eq:revision-renovacion}
\end{equation}

\begin{proposition}[Unicité formelle de la continuation normalisée]
L'équation \eqref{x_reciprocidad:eq:revision-renovacion} détermine une unique série et
\[
 \mathscr C(z)=169z^6+\frac{D_Az^7}{1-D_Az}
             =169z^6+\sum_{n=7}^{\infty}D_A^{n-6}z^n.
\]
La réalisation est analytique pour \(|D_Az|<1\).
\end{proposition}
\begin{proof}
L'élément \(1-D_Az\) a pour terme constant un et est inversible
dans l'anneau des séries formelles. Résoudre l'équation produit la série
affichée. De manière équivalente, le coefficient de degré sept vaut \(D_A\)
et chaque coefficient ultérieur vaut \(D_A\) fois le précédent. La série
géométrique converge absolument dans le disque indiqué.
\end{proof}

La normalisation du premier prolongement est une donnée opératoire
indispensable. Si l'on ne conservait que l'impulsion de degré six et
le rapport de concaténation, les séries
\[
 169z^6+\lambda\frac{D_Az^7}{1-D_Az}
\]
partageraient encore ces deux données pour tout \(\lambda\).
La règle unitaire fixe \(\lambda=1\) avant l'évaluation.
La contribution de chaque règle à l'unicité de l'opérateur de lecture
est ainsi localisée, et l'expression rationnelle conserve la continuation complète.

\subsubsection{Constante de Planck réduite et retour de la section d'action}

Dans la carte dimensionnelle de \eqref{x_reciprocidad:eq:el-secciones-accion}, la section
antérieure au retour est la construction du modèle destinée à la comparaison
avec la constante de Planck réduite :
\[
 \hbar_{\mathrm{pre}}=H_5\,10^{-34}\mathcal U_S.
\]
La section postérieure conserve également la compensation régionale
\(90\mathscr C(\alpha)/\pi\). Toutes deux appartiennent à la même droite
d'action et leur quotient \(R_{\mathrm{act}}\) mesure l'effet multiplicatif
du retour. Le passage à \(h\) par \(2\pi\) utilise la relation
entre la paramétrisation angulaire et un tour complet.

Cette construction intervient effectivement dans
\(R_{\mathrm{act}}^{80-54\alpha+6\Delta_4}\). L'électron conserve
sa coordonnée centrale et sa trajectoire ; l'échelle chronologique ne
remplace pas cette structure par un quantum indifférencié. On ne
multiplie pas non plus à nouveau par \(10^{-34}\) une masse déjà exprimée en MeV :
la décade appartient à la section d'action, se simplifie dans son quotient
sans dimension et ne réapparaît que là où agit la réalisation
dimensionnelle correspondante.

La comparaison numérique de \(H_5\), de la section retournée et de
\(h/(2\pi)\) est présentée dans la section~\ref{x_reciprocidad:sec:revision-planck-contraste}.
On distingue ainsi la dénomination de l'objet physique, l'équation que le
modèle lui attribue et la précision avec laquelle cette équation le reproduit.

% END OWNER


\subsection{Deux lecteurs du retour électronique}

La mémoire électronique est représentée dans un mot tritique de longueur \(108\) et dans un mot signé de douze phases. Pour que le calcul de l'exposant soit reproductible, nous explicitons la projection chronologique utilisée. Soient initialement \((r_0,c_0)=(5,5)\) et \((h_0,v_0)=(1,1)\). À chaque pas, le sélecteur répète \((+1,-1,0)\). Avec une orientation fixe au sein de chaque bloc de vingt-sept pas, la règle visible est
\[
 \begin{array}{c|c|c}
 \text{mode}&\text{entier lu}&\text{transport cellulaire}\\\hline
 +1&r+c&c\mapsto1+(c-1+h)\bmod9\\
 -1&rc&r\mapsto1+(r-1+v)\bmod9\\
 0&9&\text{sans déplacement}
 \end{array}
\]
Le symbole émis est \(\rho_9(\text{entier lu})\bmod3\).
Après les étapes \(27,54,81\), on inverse simultanément
\(h\) et \(v\). Cette règle calcule la projection observable ; la
mise à jour de l'état conserve en outre le résidu, le quotient,
la retenue, le feuillet, la frontière et la mémoire. Le registre de
\(108\) symboles n'est pas identifié à l'état complet.

Dans chaque bloc complet de trois étapes, on a avancé une fois dans chaque
coordonnée ; neuf de ces blocs sont nécessaires pour revenir au centre dans
la même phase. La longueur minimale de ce retour enregistré est, par
conséquent, \(27\). Si l'on oublie la phase, la cellule est déjà atteinte après
l'étape \(26\), avant le seuil neutre qui complète le bloc ; ce n'est pas
le retour typé que compte \(\mathcal K_e\). La lecture explicite
des deux premiers blocs de retour donne
\begin{equation}
 a=100000220,\qquad b=120200000,\qquad
 \gamma_e=(a^3b^3)^2.
 \label{x_reciprocidad:eq:el-palabra}
\end{equation}
Ici, les puissances désignent la concaténation de mots. En particulier,
\(\gamma_{e,t+54}=\gamma_{e,t}\). Dans le registre de neuf étapes par
phase, les orientations centrales donnent
\begin{equation}
 \tau_e=(1,1,1,-1,-1,-1,1,1,1,-1,-1,-1).
 \label{x_reciprocidad:eq:el-tau}
\end{equation}
La périodicité de ces représentations n'équivaut pas à la réinitialisation de l'histoire qui les produit.

\begin{definition}[Lecteur des discordances cycliques]
Pour un mot \(\gamma\in\mathbb F_3^{108}\), on définit
\[
 D_k(\gamma)=\#\{t\in\mathbb Z/108\mathbb Z:
                   \gamma_t\ne\gamma_{t+k}\}.
\]
Dans le domaine des mots de période divisant \(54\), le lecteur entier est
\begin{equation}
 K_D(\gamma)=\left(D_{27}+D_{36},D_1,\frac{D_4-D_3}{2}\right).
 \label{x_reciprocidad:eq:el-kd}
\end{equation}
\end{definition}

\begin{lemma}[Intégralité et évaluation centrale]
Les \(D_k\) du domaine précédent sont pairs. Pour
\eqref{x_reciprocidad:eq:el-palabra},
\[
 \begin{array}{c|rrrrr}
 k&1&3&4&27&36\\\hline
 D_k&54&56&68&48&32
 \end{array}
\]
y \(K_D(\gamma_e)=(80,54,6)\).
\end{lemma}
\begin{proof}
L'application \(t\mapsto t+54\) apparie sans points fixes les indices de discordance et conserve leur condition ; chaque cardinal est pair. Pour le mot explicite \(a^3b^3\) de longueur \(54\), le décompte des discordances avec chaque décalage est, dans le même ordre, \(27,28,34,24,16\). Doubler le mot double ces cinq cardinaux. Enfin,
\[
 (D_{27}+D_{36},D_1,(D_4-D_3)/2)
 =(48+32,54,(68-56)/2)=(80,54,6).
\]
\end{proof}

Les déplacements \(27\) et \(36\) sont les élévations
chronologiques des canaux de \(90\) et \(120\) degrés dans une carte
où neuf étapes correspondent à trente degrés. La lecture angulaire
explique cette représentation, mais les entiers se calculent sur le
mot, sans introduire de masse ni de valeur réelle recherchée.

\begin{definition}[Lecteur de fenêtres dodécaphasiques]
Pour \(\tau\in\{-1,0,1\}^{12}\), avec des indices cycliques, soient
\begin{align*}
 C_k(\tau)&=\sum_{r=0}^{11}\tau_r\tau_{r+k},\\
 A_\ell(\tau)&=\sum_{r=0}^{11}
           \left(\sum_{j=0}^{\ell-1}\tau_{r+j}\right)^2,\\
 P_6(\tau)&=\sum_{r=0}^{5}\tau_r\tau_{r+6}.
\end{align*}
On définit
\begin{equation}
 \Omega(\tau)=4A_3(\tau)-3A_4(\tau),\qquad
 K_\Omega(\tau)=(\Omega(\tau),54,P_6(\tau)).
 \label{x_reciprocidad:eq:el-komega}
\end{equation}
Le \(54\) de ce lecteur enregistre le demi-tour chronologique déclaré.
Il ne se confond pas avec le saut local de valeur quatre d'APP.
\end{definition}

\begin{lemma}[Contraste primitif et forme de corrélation]
Parmi les contrastes linéaires entiers \(uA_3+vA_4\) qui s'annulent
chaque fois que \(A_3=3a\) et \(A_4=4a\), le couple primitif, d'
orientation \(90-120\), est \((u,v)=(4,-3)\). De plus,
\begin{equation}
 \Omega=-2C_1-4C_2-6C_3.
 \label{x_reciprocidad:eq:el-omega-corr}
\end{equation}
\end{lemma}
\begin{proof}
L'annulation pour tout \(a\) exige \(3u+4v=0\). Par coprimalité, les solutions sont \((u,v)=n(4,-3)\), et l'orientation fixe le signe du générateur primitif. En développant les carrés,
\[
 A_\ell=\ell C_0+2\sum_{k=1}^{\ell-1}(\ell-k)C_k.
\]
Par conséquent,
\(A_3=3C_0+4C_1+2C_2\) et
\(A_4=4C_0+6C_1+4C_2+2C_3\). Leur combinaison \(4A_3-3A_4\)
élimine \(C_0\) et produit \eqref{x_reciprocidad:eq:el-omega-corr}.
\end{proof}

L'unicité démontrée appartient à cette classe de contrastes linéaires
entiers ; elle n'exclut pas à elle seule toute fonctionnelle non linéaire des fenêtres.
La spécification de la classe fonctionnelle évite de transformer une identité
de deux canaux en une affirmation de sélection universelle.

\begin{theorem}[Triplet de retour de la section électronique]
Les deux lecteurs appliqués aux registres de \(\Gamma_e\) coïncident et donnent
\begin{equation}
 K_D(\gamma_e)=K_\Omega(\tau_e)=(80,54,6).
 \label{x_reciprocidad:eq:el-triple}
\end{equation}
Son évaluation aux coordonnées déjà générées est
\begin{equation}
 \kappa_e=80-54\alpha+6\Delta_4.
 \label{x_reciprocidad:eq:el-kappa}
\end{equation}
\end{theorem}
\begin{proof}
Le premier lecteur a été calculé plus haut. Pour \(\tau_e\), le produit
antipodal vaut toujours un, d'où \(P_6=6\). Le dénombrement cyclique donne
\(C_1=4\), \(C_2=-4\), \(C_3=-12\) ; donc
\[
 \Omega=-2\cdot4-4(-4)-6(-12)=80.
\]
De manière équivalente, \(A_3=44\), \(A_4=32\) et \(4\cdot44-3\cdot32=80\). Les deux triplets sont donc \((80,54,6)\). L'évaluation de la fonctionnelle \((A,B,C)\mapsto A-\alpha B+\Delta_4C\) prouve la dernière formule.
\end{proof}

Cette coïncidence correspond à la section centrale. Elle n'identifie pas les deux lecteurs dans tous les domaines de routes : l'un compare des symboles tritiques et l'autre compare des fenêtres d'un registre signé. De même, \(80\) ne s'obtient pas à partir de la période d'un mot attribué au nombre d'Euler. Son origine dans cette composition est le triplet de la route électronique.

\subsection{Composition électronique complète et normalisation}

\begin{definition}[Retour multiplicatif]
L'action de la mémoire de retour sur la coordonnée basale est
\begin{equation}
 \varepsilon_e^\beta
 =\varepsilon_e^{(0)}\exp\!\left(\kappa_e\log R_{\mathrm{act}}\right).
 \label{x_reciprocidad:eq:el-beta}
\end{equation}
\end{definition}

\begin{theorem}[Factorisation électronique positive]
Avec les lecteurs, l'orientation et la normalisation déclarés,
\eqref{x_reciprocidad:eq:el-beta} coïncide avec \eqref{x_reciprocidad:eq:el-formula-completa}. Tous
ses facteurs sont adimensionnels et la coordonnée résultante est positive.
\end{theorem}
\begin{proof}
La norme d'incompatibilité est \eqref{x_reciprocidad:eq:el-prefactor} ; le transfert modal,
le déficit et la mémoire retenue composent \eqref{x_reciprocidad:eq:el-basal}. Le
triplet \eqref{x_reciprocidad:eq:el-triple} produit \eqref{x_reciprocidad:eq:el-kappa}, et le rapport des
deux sections d'action est \eqref{x_reciprocidad:eq:el-ratio}. En substituant ces
trois identités dans \eqref{x_reciprocidad:eq:el-beta}, on obtient exactement
\eqref{x_reciprocidad:eq:el-formula-completa}. L'étape basale est positive, tout comme
l'exponentielle d'un nombre réel. Le facteur commun de dimension d'
action s'est déjà simplifié lors de la formation de \(R_{\mathrm{act}}\).
\end{proof}

L'évaluation de cette composition commence par
\begin{equation}
 \varepsilon_e^\beta
 =0.5109989506900008086082571648\ldots.
 \label{x_reciprocidad:eq:el-beta-decimal}
\end{equation}
Les chiffres sont une évaluation ultérieure de la formule, et non des données utilisées
pour sélectionner \(22\), \(15\) ou \((80,54,6)\). La représentation
numérique du terme basal par un préfixe fini introduit uniquement son
erreur de troncature ; elle ne modifie pas la formule exacte qui le définit.

L'incidence d'une normalisation torsionnelle différente peut être localisée sans la comparer à une mesure. Écrivons \(B=\frac{\sqrt3}{4}[\exp(\varphi/\pi^2)-22\alpha^3]\), \(R=R_{\mathrm{act}}\) et
\[
 \mathcal E(d)=B(1+15d)R^{80-54\alpha+6d}.
\]
Pour les deux défauts de \eqref{x_reciprocidad:eq:el-delta},
\begin{equation}
 \frac{\mathcal E(\Delta_4)}{\mathcal E(d_4)}
 =\frac{1+15\Delta_4}{1+15d_4}
   R^{6(\Delta_4-d_4)}>1.
 \label{x_reciprocidad:eq:el-normalizacion-beta}
\end{equation}
La première inégalité torsionnelle et \(R>1\) prouvent qu'il ne s'agit pas de la même coordonnée. Si l'on ne change que le défaut de l'exposant, en conservant le terme basal complet, le rapport des deux sorties est \(R^{6(\Delta_4-d_4)}\), lui aussi différent de un. Une très petite différence reste une différence de formule ; elle ne peut pas être attribuée à une simple conversion d'unités.

\subsection{Réalisation dimensionnelle : énergie, masse et échelle chronologique}

La coordonnée électronique agit finalement sur une droite dimensionnelle.
Soit \(\mathcal L_E\) une droite d'énergie de base positive
\(\mathcal U_E\). La réalisation est
\begin{equation}
 E_e^{(0)}=\varepsilon_e^{(0)}\mathcal U_E,\qquad
 E_e^\beta=\varepsilon_e^\beta\mathcal U_E,
 \qquad m_e^\beta=\frac{E_e^\beta}{c_{\mathrm{int}}^2}.
 \label{x_reciprocidad:eq:el-dimensional}
\end{equation}
La dernière expression appartient à la droite de masse lorsque
\(c_{\mathrm{int}}\) porte la dimension d'une vitesse dans la réalisation
déclarée. Dans la carte comparative \(\mathcal U_E=1\,\mathrm{MeV}\),
\eqref{x_reciprocidad:eq:el-beta-decimal} exprime une énergie en MeV ; la masse
correspondante s'exprime en \(\mathrm{MeV}/c^2\). Il est fréquent
d'abréger les deux grandeurs en prenant \(c=1\), mais cette abréviation ne
doit pas masquer l'application dimensionnelle.

L'échelle d'action admet aussi une réalisation action--horloge. Si
\(t_0\) est une section temporelle positive, l'horloge de \(108\) étapes
définit
\begin{equation}
 \omega_0=\frac{2\pi}{108t_0},\qquad
 E_{\mathrm{clk}}=\hbar_{\mathrm{ret}}\omega_0,\qquad
 m_{\mathrm{clk}}=E_{\mathrm{clk}}c_{\mathrm{int}}^{-2}.
 \label{x_reciprocidad:eq:el-cuanto-reloj}
\end{equation}
L'action multipliée par la fréquence a la dimension d'une énergie. Ce quantum
chronologique est une section dimensionnelle commune ; il n'est pas, par définition,
la section électronique. Le mot de route et ses facteurs ne s'
obtiennent pas à partir de l'uniformité de l'horloge.

\begin{proposition}[Annulation de l'échelle commune]
Le facteur \(10^{-34}\mathcal U_S\) de \eqref{x_reciprocidad:eq:el-secciones-accion} ne fait pas partie de \(\kappa_e\). Sur une même droite d'énergie,
\begin{equation}
 \frac{E_e^\beta}{E_e^{(0)}}=R_{\mathrm{act}}^{\kappa_e}.
 \label{x_reciprocidad:eq:el-razon-energia}
\end{equation}
Le résultat est indépendant du choix de la base dimensionnelle
commune.
\end{proposition}
\begin{proof}
Lorsqu'on forme \(\hbar_{\mathrm{pre}}/\hbar_{\mathrm{ret}}\), le facteur
\(10^{-34}\mathcal U_S\) se simplifie exactement. Lorsqu'on divise les deux
expressions d'énergie de \eqref{x_reciprocidad:eq:el-dimensional}, \(\mathcal U_E\) se simplifie
également, et \eqref{x_reciprocidad:eq:el-beta} donne le rapport indiqué.
\end{proof}

Une coordonnée déjà exprimée en MeV n'est donc pas multipliée une nouvelle fois par \(10^{-34}\). Si l'on choisit \(E_{\mathrm{clk}}\) comme base de \(\mathcal L_E\), les coordonnées doivent être transportées par le facteur de changement de base ; conserver sans changement le nombre de la carte d'un MeV reviendrait à identifier deux bases sans le démontrer. C'est le point précis où l'échelle d'action entre dans la réalisation absolue et cesse d'intervenir dans le correcteur relatif.

Il existe également une réalisation linéaire qui rend visible la séparation
entre horloge et généalogie. Dans
\(\ell^2(\mathbb Z/108\mathbb Z;\mathbb R)\), le décalage
cyclique \(S\) a pour espace fixe
\[
 \ker(S-I)=\mathbb R n_0,\qquad
 n_0=\frac1{\sqrt{108}}\sum_{t=0}^{107}\delta_t.
\]
Si \(g_{\Gamma_e}\) est le symbole de la route dans la linéarisation libre
des généalogies et \(\mathcal U_M\) une base de la droite de masse,
on définit l'application de rang un
\begin{equation}
 \mathfrak C_e:
 \mathbb R n_0\otimes\mathbb R g_{\Gamma_e}\longrightarrow\mathcal L_M,
 \qquad
 \mathfrak C_e(n_0\otimes g_{\Gamma_e})
 =\varepsilon_e^\beta\mathcal U_M.
 \label{x_reciprocidad:eq:el-rango-uno}
\end{equation}
L'application réunit la factorisation précédente une fois ses bases déclarées.
Le vecteur uniforme \(n_0\) ne sélectionne pas \(\varepsilon_e^\beta\) :
cette coordonnée procède d'APP, du retour orienté et de ses lecteurs.
On distingue ainsi une réalisation algébrique du mode nul d'une
identification de l'électron au quantum de l'horloge.

La conclusion interne est la factorisation positive d'une section
centrale, avec ses opérations, sa mémoire et ses normalisations explicites.
L'attribution à un électron physique ajoute la représentation de charge, le
spin, le couplage et la carte d'unités correspondante ; elle ne se
déduit pas exclusivement d'une coïncidence décimale. La comparaison
métrologique est présentée après le prolongement exceptionnel et ne modifie
les entrées d'aucune des constructions réalisées ici.


% BEGIN OWNER /Users/ruben/Documents/New project/output/EXTREMA_MEDIA_RAZON_RECIPROCIDAD_ES_EN_20260918/ARTICULO/ES/source/sections/revision_electron_argumento.tex
\subsubsection{Dépendance structurelle de la coordonnée électronique}

L'équation électronique exprime une composition hiérarchique. Son premier
facteur, \(\sqrt3/4\), appartient à la structure matricielle centrale et fixe
une normalisation géométrique. L'exposant \(\varphi/\pi^2\) utilise les
coordonnées régionales déjà engendrées dans l'orientation indiquée de l'opérateur
de lecture surfacique--volumique. Le terme \(22\alpha^3\) incorpore le couplage
dodécaphasique à la correction cubique. Le facteur \(1+15\Delta_4\) conserve
la normalisation torsionnelle globale. Enfin, le quotient des sections
d'action agit avec l'exposant déterminé par les opérateurs de lecture du retour.
Chaque opération possède ainsi un antécédent distinct au sein de la même
construction.

L'articulation devient plus précise lorsqu'on considère ses dépendances.
La norme centrale peut être calculée avant l'évaluation des coordonnées
régionales. L'exponentielle et la correction cubique requièrent ces
coordonnées et la clôture de \(\alpha\). La mémoire multiplicative requiert
en outre les sections d'action et la trajectoire de \(108\) positions.
Connaître l'expression basale n'équivaut donc pas à avoir évalué la
coordonnée complète. La succession d'opérations n'ajoute pas une masse
externe au point \((5,5)\) ; elle développe les opérateurs de lecture attribués à sa
section persistante.

La distinction entre échelle et structure fournit également une
interprétation de la comparaison physique. La masse est la réalisation
dimensionnelle de la coordonnée complète, tandis que le centre, l'involution
et la lacune décrivent les relations qui la produisent. Un changement commun
d'unité modifie les coordonnées dimensionnelles, mais laisse inchangés
les rapports sans dimension et les identités de transport démontrées.
C'est en ce sens que l'électron conserve son caractère structurel
lorsqu'il s'exprime par rapport à une échelle chronologique.

% END OWNER



% BEGIN OWNER /Users/ruben/Documents/New project/output/EXTREMA_MEDIA_RAZON_RECIPROCIDAD_ES_EN_20260918/ARTICULO/ES/source/sections/realizacion_electromagnetica.tex
\newpage
\subsection{Transport lorentzien et représentation électromagnétique}

La section électronique fournit une structure et une coordonnée
avant son application à une interaction. La réalisation lorentzienne
de l'incidence cubique détermine, avec la convention temporelle déclarée,
l'espace des formes dans lequel cette interaction peut s'exprimer.
Pour la formuler, on fixe une réalisation locale orientée \((\mathcal M,g)\)
de dimension quatre, avec des fibres tangentes identifiées au modèle
\(Q_\square\). On spécifie en outre une connexion \(\nabla\) avec
\(\nabla g=0\), une un-forme potentielle \(A\in\Omega^1(\mathcal M)\),
sa courbure \(F=dA\) et une représentation de charge de la section.
Ces données sont les arguments de la représentation physique qui suit.

La dualité de Hodge satisfait \(\star^2=-I\) sur les deux-formes. La séparation électrique et magnétique correspond à un choix temporel dans la décomposition \(1+3\) ; \(F\) et \(\star F\) conservent leur relation au sein du même espace de deux-formes. La neutralité d'une représentation de charge signifie \(q=0\), tandis que la section peut conserver orientation, phase et transport interne. L'absence de couplage électrique dans cette représentation a un contenu distinct de l'absence d'information de phase.

Dans cette réalisation dimensionnelle, soient \(m_e>0\), \(q\) la charge,
\(\gamma(s)\) une trajectoire paramétrée et \(u=\dot\gamma\)
sa quadrivitesse. La force de Lorentz s'exprime par
\begin{equation}
 m_e\nabla_u u^\mu=qF^\mu{}_{\nu}u^\nu.
 \label{x_reciprocidad:eq:lorentz-electron}
\end{equation}
L'antisymétrie de \(F\) implique
\(F_{\mu\nu}u^\mu u^\nu=0\). Par conséquent,
\[
 \frac{d}{ds}g(u,u)=2g(\nabla_u u,u)=0.
\]
La normalisation de la trajectoire est conservée dès lors que \(\nabla\)
est métrique et que \eqref{x_reciprocidad:eq:lorentz-electron} est satisfaite. La masse utilisée
est la réalisation dimensionnelle de la section électronique précédente ;
le paramètre de couplage est identifié dans la convention électromagnétique
adoptée. La déduction concrète du potentiel et d'une configuration
de champ constitue un problème supplémentaire assorti de conditions de frontière
propres. Le résultat exposé ici détermine le domaine, les données
de la représentation et la conservation qu'implique son équation.

L'inversion des routes du TPK et l'orientation d'une ligne électronique
offrent alors une comparaison précise : chacune possède son involution,
son transport et sa représentation. L'antécédent de Wheeler et Feynman
situe historiquement la question ; l'équivalence de deux réalisations
se décide au moyen de leurs applications et conditions de compatibilité.
Cette disposition permet de maintenir le lien entre l'origine arithmétique de l'électron
et la formulation de son interaction, en explicitant la fonction
de chaque étape.

% END OWNER

